\DocumentMetadata{
  pdfversion=2.0,
  pdfstandard=ua-2,
  lang=en-US,
  tagging=on,
  tagging-setup={math/setup=mathml-SE,tabsorder=structure}
}
\documentclass[11pt,letterpaper]{article}
\usepackage[margin=0.68in,includefoot]{geometry}
\usepackage{newcomputermodern}
\usepackage{amsmath,amscd,mathtools}
\usepackage{tikz,adjustbox}
\providecommand{\square}{\mdlgwhtsquare}
\usepackage[hidelinks]{hyperref}
\hypersetup{
  pdftitle={Algebra PhD Exam, September 1999},
  pdfauthor={University of Florida Department of Mathematics},
  pdfsubject={Graduate mathematics examination},
  pdfdisplaydoctitle=true,
  bookmarksopen=true
}
\setcounter{secnumdepth}{0}
\setlength{\parindent}{0pt}
\setlength{\parskip}{0.28em}
\setlength{\emergencystretch}{3em}
\setlength{\leftmargini}{2.15em}
\setlength{\leftmarginii}{2.65em}
\setlength{\leftmarginiii}{3em}
\renewcommand{\labelenumi}{\textbf{\theenumi.}}
\pagestyle{empty}
\raggedbottom
\allowdisplaybreaks
\newenvironment{examproblems}[1]{%
  \begin{enumerate}%
  \setcounter{enumi}{#1}%
  \setlength{\topsep}{0.35em}%
  \setlength{\partopsep}{0pt}%
  \setlength{\itemsep}{0.48em}%
  \setlength{\parsep}{0.18em}%
}{\end{enumerate}}
\newenvironment{examsubparts}{%
  \begin{enumerate}%
  \setlength{\topsep}{0.18em}%
  \setlength{\partopsep}{0pt}%
  \setlength{\itemsep}{0.16em}%
  \setlength{\parsep}{0.1em}%
}{\end{enumerate}}
\newenvironment{examinstructions}{%
  \par\begingroup\itshape
}{%
  \par\endgroup\vspace{0.18em}
}
\makeatletter
\renewcommand\section{\@startsection{section}{1}{0pt}%
  {-1.25ex plus -.25ex minus -.1ex}%
  {0.55ex plus .1ex}%
  {\normalfont\large\bfseries}}
\renewcommand{\@maketitle}{%
  \newpage\null\vskip -1em%
  {\footnotesize\bfseries
    \href{https://www.ufl.edu/}{University of Florida}, \href{https://math.ufl.edu/}{Department of Mathematics}\par}%
  \vskip -0.5em%
  \tagstructbegin{tag=Title}%
    {\LARGE\bfseries\href{https://gma.math.ufl.edu/exams/algebra-phd/}{Algebra PhD Exam}, September 1999\par}%
  \tagstructend%
  \vskip 0.25em%
  \tagmcbegin{artifact}%
    \hrule height 1pt%
  \tagmcend%
  \vskip 1em%
}
\makeatother
\title{Algebra PhD Exam, September 1999}
\author{}
\date{}
\begin{document}
\maketitle
\thispagestyle{empty}
\begin{examinstructions}
Please do 7 out of the 11 problems below.
\end{examinstructions}
\begin{examproblems}{0}
\item
Here \(\mathbb{Z}_n\) denotes the cyclic group of order~\(n\).
\begin{examsubparts}
\item[\textnormal{(a).}]
Let~\(p\) be a prime. How many subgroups of order \(p^2\) does the abelian group \(\mathbb{Z}_{p^3}\oplus\mathbb{Z}_{p^2}\) have?
\item[\textnormal{(b).}]
What are the elementary divisors of the group \(\mathbb{Z}_{26}\oplus\mathbb{Z}_{42}\oplus\mathbb{Z}_{49}\oplus\mathbb{Z}_{200}\)? What are its invariant factors?
\item[\textnormal{(c).}]
Determine up to isomorphism all abelian groups of order 20.
\end{examsubparts}
\item
This problem has two parts.
\begin{examsubparts}
\item[\textnormal{(a).}]
Exhibit an automorphism of the cyclic group of order 6 that is not an inner automorphism.
\item[\textnormal{(b).}]
Show that if a normal subgroup~\(N\) of order~\(p\) (\(p\) prime) is contained in a group~\(G\) of order \(p^n\), then~\(N\) is in the center of~\(G\).
\end{examsubparts}
\item
This problem has three parts.
\begin{examsubparts}
\item[\textnormal{(a).}]
Is the additive group~\(\mathbb{Q}\) of rationals a divisible abelian group?
\item[\textnormal{(b).}]
Show that no nonzero finite abelian group is divisible.
\item[\textnormal{(c).}]
Show that no nonzero free abelian group is divisible.
\end{examsubparts}
\item
Let~\(R\) and~\(S\) be rings and \(A_R, {}_R B_S, C_S\) (bi)-modules. Demonstrate an isomorphism of abelian groups between \(\operatorname{Hom}_S(A\otimes_R B,C)\) and \(\operatorname{Hom}_R(A,\operatorname{Hom}_S(B,C))\). (You should state very clearly what must be checked.)
\item
Let~\(R\) be a ring, let 
\[
0\longrightarrow A\xrightarrow{f}B\xrightarrow{g}C\longrightarrow 0
\]
 be a short exact sequence of left~\(R\)-modules and let~\(D\) be a right~\(R\)-module. Show that 
\[
0\longrightarrow D\otimes_R A\xrightarrow{1_D\otimes f}D\otimes_R B\xrightarrow{1_D\otimes g}D\otimes_R C\longrightarrow 0
\]
 is a short exact sequence of abelian groups under each one of the following hypotheses:
\begin{examsubparts}
\item[\textnormal{(a).}]
\(0\longrightarrow A\xrightarrow{f}B\xrightarrow{g}C\longrightarrow 0\) is split exact.
\item[\textnormal{(b).}]
\(R\) has an identity and~\(D\) is a free right~\(R\)-module.
\item[\textnormal{(c).}]
\(R\) has an identity and~\(D\) is a projective unitary right~\(R\)-module.
\end{examsubparts}
\item
Let~\(A\) be a commutative ring with 1 such that every element \(x\in A\) satisfies \(x^n=x\) for some \(n>1\), depending on~\(x\). Prove that
\begin{examsubparts}
\item[\textnormal{(a).}]
every prime ideal of~\(A\) is maximal.
\item[\textnormal{(b).}]
the Jacobson radical of~\(A\) is \((0)\).
\end{examsubparts}
\item
Let~\(C\) and~\(D\) be categories, and let~\(S\) and~\(T\) be covariant functors from~\(C\) to~\(D\).
\begin{examsubparts}
\item[\textnormal{(a).}]
Define a natural isomorphism \(\alpha:S\to T\).
\item[\textnormal{(b).}]
If~\(B\) is a unitary left module over a ring~\(R\) with identity, show that there is a “natural” isomorphism of modules 
\[
\alpha_B:R\otimes_R B\cong B.
\]
\end{examsubparts}
\item
This problem has two parts.
\begin{examsubparts}
\item[\textnormal{(a).}]
Give an example of commutative rings \(R,A\) with \(R\subset A\), where~\(A\) is Noetherian but~\(R\) is not.
\item[\textnormal{(b).}]
Let \(R=K[x,y]\) be the polynomial ring in two variables over the field~\(K\), let~\(A\) denote the ideal \((xy,y^2)\) in~\(R\) and let \(c\in K\). Show that 
\[
A=(y)\cap(x+cy,y^2)
\]
 and that this is a primary decomposition.
\end{examsubparts}
\item
This problem has two parts.
\begin{examsubparts}
\item[\textnormal{(a).}]
Show that if~\(S\) is a multiplicative subset of a Dedekind domain~\(R\) (with \(1\in S\), \(0\notin S\)), then \(S^{-1}R\) is a Dedekind domain.
\item[\textnormal{(b).}]
Show that if~\(I\) is a nonzero ideal in a Dedekind domain, then \(R/I\) is an Artinian ring.
\end{examsubparts}
\item
Let~\(p\) be a prime and \(n\geq 1\) an integer. Show that~\(F\) is a finite field with \(p^n\) elements if and only if~\(F\) is a splitting field of \((x^{p^n}-x)\) over the field \(\mathbb{Z}_p\) of~\(p\) elements.
\item
This problem has two parts.
\begin{examsubparts}
\item[\textnormal{(a).}]
Give an example of fields \(M\supset L\supset K\) such that \(M/L\) and \(L/K\) are normal extensions, but \(M/K\) is not normal. Justify your answer.
\item[\textnormal{(b).}]
Let \(L\supset K\) be fields such that \([L:K]=8\). Prove that there exist \(\alpha,\beta,\gamma\in L\) such that \(L=K(\alpha,\beta,\gamma)\).
\end{examsubparts}
\end{examproblems}
\end{document}
